\begin{slide}
  \begin{block}{Problem}
    Given a multiplication $f_1 \cdot f_2 = p$ of unknown base $b$, what could a possible $b$ be?
  \end{block}
  \pause
  \begin{block}{Insight}
    \begin{itemize}
      \item If $f_1$ and $f_2$ are multiplied in a base greater than $b$, the resulting product
        is always smaller than $p$ in this base.
      \item If $f_1$ and $f_2$ are multiplied in a base smaller than $b$, the resulting product
        is always greater than $p$ in this base.
      \item Since all digits are smaller than or equal to $2^{30}$ and we have at most $1\,000$
        digits, $b$ must be smaller than $2^{61}$.
    \end{itemize}
  \end{block}
\end{slide}
\begin{slide}
  \begin{block}{Lemma}
    Let $f_1$ and $f_2$ be factors in a multiplication with no more than $n$ digits, $X$ the
    highest occurring digit and $p$ the product with no more than $n$ digits. If the
    multiplication is performed in a sufficiently high base such that no carries occur and
    $X > (n + 1)^{2}$, the smallest digit of $p$ that is greater than $X$ is smaller
    than $2 \cdot X^{2}$.
  \end{block}
  \pause
  \begin{block}{Proof}
    \begin{itemize}
      \item Let $a_i$ be the digits of $f_1$, $b_i$ the digits of $f_2$ and $c_i$ the digits of $p$.
      \item Since no carries occurred, every digit of $p$ can be calculated with
        $c_k = \sum_{i + j = k} a_i \cdot b_j$. The sum for $c_k$ has at most $k + 1$ summands.
      \item Let's assume $c_{k_2}$ with $c_{k_2} \ge 2 \cdot X^{2}$ is the smallest digit greater than $X$.
    \end{itemize}
  \end{block}
\end{slide}
\begin{slide}
  \begin{block}{Proof cont.}
    \begin{itemize}
      \item There must be at least one pair $a_{i_1}$, $b_{j_1}$ in the sum to calculate $c_{k_2}$ such that
        $\min(a_{i_1}, b_{j_1}) \ge \frac{X}{k_2 + 1}$. Otherwise $c_{k_2} < X^2$.
      \item Among the other $k_2$ summands, there must be a second pair $a_{i_2}$, $b_{j_2}$ with
        $\min(a_{i_2}, b_{j_2}) \ge \frac{X}{k_2}$. Otherwise $c_{k_2} < X^2 + a_{i_1} b_{j_1} \le 2 \cdot X^2$.
      \item Without limiting the generality, let $i_1 < i_2$ and $j_2 < j_1$. Then in the sum of an earlier product digit
        $a_{i_1} \cdot b_{j_2}$ must have been a summand. Let this digit be $c_{k_1}$.
      \item Then $c_{k_1} \ge a_{i_1} \cdot b_{j_2} \ge \frac{X}{k_2 + 1} \cdot \frac{X}{k_2} >
        \frac{X}{n + 1} \cdot \frac{X}{n + 1} = \frac{X^2}{(n + 1)^2}$
      \item Since our lemma assumes $X > (n + 1)^2$, $c_{k_2} > c_{k_1} > X$ $\Rightarrow$ $c_{k_2}$ is not the smallest
        digit greater than $X$ $\Rightarrow$ contradiction to assumption.
    \end{itemize}
  \end{block}
\end{slide}
\begin{slide}
  \begin{block}{Lemma}
    Let $f_1$ and $f_2$ be factors in a multiplication with no more than $n$ digits, $X$ the
    highest occurring digit and $p$ the product with no more than $n$ digits. If the
    multiplication is performed in a sufficiently high base such that no carries occur and
    $X > (n + 1)^{2}$, the smallest digit of $p$ that is greater than $X$ is smaller
    than $2 \cdot X^{2}$.
  \end{block}
  \begin{block}{Conclusion}
    \begin{itemize}
      \item Since in our problem $X = 2^{30} > (1\,000 + 1)^2$, we can use this lemma to determine
        the highest possible base.
      \item In our problem no digit in $p$ is greater than $2^{30}$, so the highest possible base must
        be slightly smaller than if we had no carries.
      \item $\Rightarrow$ highest possible base is smaller than $2 \cdot (2^{30})^2 = 2^{61}$
    \end{itemize}
  \end{block}
\end{slide}
\begin{slide}
  \begin{block}{Problem}
    Given a multiplication $f_1 \cdot f_2 = p$ of unknown base $b$, what could a possible $b$ be?
  \end{block}
  \begin{block}{Solution}
    \begin{itemize} 
      \item Write a multiplication routine that can multiply two numbers in a given base.
      \item Use binary search to determine the correct base.
      \pause
      \item Factorisation approach only leads to accepted solution if a fast factorisation algorithm is used.
    \end{itemize}
  \end{block}
\end{slide}
